\begin{summarybox}
	\textbf{\textcolor{CSummary}{一句话核心}}

	\textbf{弦切角等于它所夹弧所对的圆周角}。

	\medskip
	\textbf{\textcolor{CSummary}{使用条件}}
	\begin{itemize}[leftmargin=*, itemsep=0.5em, topsep=0.4em]
		\item \textbf{条件1（有切线与弦）：} 图形中存在切线和过切点的弦，构成弦切角。
		\item \textbf{条件2（有对应圆周角）：} 圆上有异于切点的点，其角对同一弧。
	\end{itemize}

	\medskip
	\textbf{\textcolor{CSummary}{关键提醒}}
	\begin{itemize}[leftmargin=*, itemsep=0.5em, topsep=0.4em]
		\item \textbf{易错点（弧的对应）：} 必须确保弦切角所夹弧与圆周角所对弧一致。
		\item \textbf{检查项（角顶点）：} 弦切角顶点必须为切点。
	\end{itemize}
\end{summarybox}
